% Options for packages loaded elsewhere
\PassOptionsToPackage{unicode}{hyperref}
\PassOptionsToPackage{hyphens}{url}
\documentclass[
  11pt,
]{article}
\usepackage{xcolor}
\usepackage[margin=1in]{geometry}
\usepackage{amsmath,amssymb}
\setcounter{secnumdepth}{-\maxdimen} % remove section numbering
\usepackage{iftex}
\ifPDFTeX
  \usepackage[T1]{fontenc}
  \usepackage[utf8]{inputenc}
  \usepackage{textcomp} % provide euro and other symbols
\else % if luatex or xetex
  \usepackage{unicode-math} % this also loads fontspec
  \defaultfontfeatures{Scale=MatchLowercase}
  \defaultfontfeatures[\rmfamily]{Ligatures=TeX,Scale=1}
\fi
\usepackage{lmodern}
\ifPDFTeX\else
  % xetex/luatex font selection
\fi
% Use upquote if available, for straight quotes in verbatim environments
\IfFileExists{upquote.sty}{\usepackage{upquote}}{}
\IfFileExists{microtype.sty}{% use microtype if available
  \usepackage[]{microtype}
  \UseMicrotypeSet[protrusion]{basicmath} % disable protrusion for tt fonts
}{}
\makeatletter
\@ifundefined{KOMAClassName}{% if non-KOMA class
  \IfFileExists{parskip.sty}{%
    \usepackage{parskip}
  }{% else
    \setlength{\parindent}{0pt}
    \setlength{\parskip}{6pt plus 2pt minus 1pt}}
}{% if KOMA class
  \KOMAoptions{parskip=half}}
\makeatother
\usepackage{graphicx}
\makeatletter
\newsavebox\pandoc@box
\newcommand*\pandocbounded[1]{% scales image to fit in text height/width
  \sbox\pandoc@box{#1}%
  \Gscale@div\@tempa{\textheight}{\dimexpr\ht\pandoc@box+\dp\pandoc@box\relax}%
  \Gscale@div\@tempb{\linewidth}{\wd\pandoc@box}%
  \ifdim\@tempb\p@<\@tempa\p@\let\@tempa\@tempb\fi% select the smaller of both
  \ifdim\@tempa\p@<\p@\scalebox{\@tempa}{\usebox\pandoc@box}%
  \else\usebox{\pandoc@box}%
  \fi%
}
% Set default figure placement to htbp
\def\fps@figure{htbp}
\makeatother
\setlength{\emergencystretch}{3em} % prevent overfull lines
\providecommand{\tightlist}{%
  \setlength{\itemsep}{0pt}\setlength{\parskip}{0pt}}
\usepackage{amsmath}
\usepackage{amssymb}
\usepackage{booktabs}
\usepackage{longtable}
\usepackage{array}
\usepackage{float}
\usepackage{setspace}
\usepackage{graphicx}
\usepackage{caption}
\usepackage{hyperref}
\usepackage{bookmark}
\IfFileExists{xurl.sty}{\usepackage{xurl}}{} % add URL line breaks if available
\urlstyle{same}
\hypersetup{
  pdftitle={Reel Intelligence: MovieLens Clustering},
  pdfauthor={Ali Jabir; Isis Martínez; Susan Peck},
  hidelinks,
  pdfcreator={LaTeX via pandoc}}

\title{Reel Intelligence: MovieLens Clustering}
\usepackage{etoolbox}
\makeatletter
\providecommand{\subtitle}[1]{% add subtitle to \maketitle
  \apptocmd{\@title}{\par {\large #1 \par}}{}{}
}
\makeatother
\subtitle{Final Project Report, Spring 2026}
\author{Ali Jabir \and Isis Martínez \and Susan Peck}
\date{May 2, 2026}

\begin{document}
\maketitle

\begin{center}
\textbf{Clustering Analysis (Math 252)}\\
Professor: Dr. Cristina Tortora\\
Department of Computer Science\\
San Jose State University
\end{center}

\section{Abstract}\label{abstract}

TBD

{[}NOTE: ADD ANY OTHER INFORMATION NEEDED WITH CODE AND RESULTS LATER{]}

{[}ADD 3 DIFFERENT DATASETS{]}

\section{Introduction}\label{introduction}

Clustering is a core unsupervised learning technique used to identify
structure in high-dimensional data without labeled outcomes. In
movie-rating data, clustering offers a different perspective from
prediction-based recommender systems by focusing on groups of users with
similar behavior rather than forecasting future ratings.

Although this study is motivated by data commonly used in recommender
systems, it does not evaluate recommendation accuracy or predictive
performance. Instead, it assesses clustering quality in learned
embedding spaces. Coherent cluster structure in these embeddings may
still have downstream relevance for personalization, recommendation, and
content organization.

The MovieLens datasets are widely used benchmarks in machine learning
and recommender-systems research because they contain rich user-movie
rating information. In this project, we use the MovieLens 1M dataset and
cluster user embeddings learned from the rating data. We consider
embedding dimensions of 64, 128, and 256 and compare four clustering
methods: k-means, probabilistic distance (PD) clustering, reduced
k-means, and factorial probabilistic distance clustering (FPDC). The
goal is to evaluate how embedding dimension and clustering method
jointly affect the quality, stability, and interpretability of
clustering solutions.

This report is organized as follows. \hyperref[background]{Section II}
reviews the clustering methods used in the study.
\hyperref[dataset]{Section III} describes the dataset and preprocessing
pipeline. \hyperref[methods-and-experimental-setup]{Section IV} presents
the embedding model and clustering implementation.
\hyperref[evaluation-framework]{Section V} defines the evaluation
framework. \hyperref[results]{Section VI} reports the experimental
results. \hyperref[conclusion-and-future-work]{Section VII} concludes
and discusses future work.

\section{Background}\label{background}

Clustering partitions observations into groups according to a notion of
similarity or distance. In this study, clustering is performed in
learned embedding spaces, so the geometry of the representation plays a
central role in determining cluster quality.

We consider four partition-based methods that differ along two key
dimensions: (i) \textbf{hard vs.~probabilistic assignment}, and (ii)
\textbf{clustering in the full space vs.~a learned lower-dimensional
subspace}.

\textbf{k-means} provides a natural baseline. It assigns each
observation to a single cluster by minimizing within-cluster variation,
and is computationally efficient and widely used. However, it implicitly
favors roughly spherical, equally sized clusters and can be sensitive to
initialization, scaling, and outliers.

\textbf{Probabilistic distance (PD) clustering} replaces hard
assignments with probabilistic memberships determined by distances to
cluster centers. This produces a soft partition in which observations
can belong to multiple clusters to varying degrees. As a result, PD
clustering is typically more robust to overlap and noise, but at
increased computational cost relative to k-means.

High-dimensional settings motivate methods that jointly perform
clustering and dimension reduction. In such cases, cluster structure may
be more clearly expressed in a lower-dimensional subspace than in the
original feature space.

\textbf{Reduced k-means} extends k-means by simultaneously estimating a
low-dimensional projection and a partition within that subspace. By
concentrating on directions that are most relevant for clustering, it
can recover structure that is obscured in the full space, while
retaining the simplicity of hard assignments.

\textbf{Factorial probabilistic distance clustering (FPDC)} combines the
ideas of PD clustering and dimension reduction. It estimates a latent
subspace together with probabilistic cluster memberships, effectively
performing soft clustering in a learned representation. This allows FPDC
to accommodate overlapping clusters while also mitigating the effects of
high dimensionality.

From this perspective, the four methods can be viewed along two
orthogonal axes: assignment (hard vs.~probabilistic) and representation
(original space vs.~learned subspace). These differences suggest
potential trade-offs in robustness, interpretability, and computational
cost.

In addition to clustering quality, computational efficiency is an
important consideration in high-dimensional settings. Methods that
reduce dimension can lower per-iteration cost and memory usage, but may
introduce additional overhead from estimating the subspace. For this
reason, this study evaluates both clustering performance and
computational scalability.

Because clustering is unsupervised, validation relies on internal
criteria, stability, and interpretability rather than labeled outcomes.
In particular, the methods considered here may differ in their
sensitivity to sampling variability. Probabilistic approaches can
accommodate ambiguity in cluster boundaries, while subspace-based
methods depend on how well the learned representation captures the
underlying structure.

To assess these differences, this study evaluates clustering
\textbf{stability under bootstrap resampling}, using agreement between
partitions as a primary criterion. Stability provides a complementary
perspective to internal fit measures, especially in settings where
cluster structure may be weak or sensitive to representation.
Accordingly, the evaluation considers multiple complementary metrics
rather than relying on a single measure.

\section{Dataset}\label{dataset}

We use the MovieLens 1M dataset as the primary dataset for this project.
MovieLens 1M is a standard benchmark in recommender-systems research and
contains approximately one million movie ratings from roughly six
thousand users on roughly four thousand movies. Its size makes it
suitable for experimenting with learned embeddings and multiple
clustering procedures while remaining computationally manageable.

In addition to MovieLens 1M, we use sampled subsets of the MovieLens 20M
dataset for benchmarking. These subsets are used to assess whether
clustering results remain stable across different samples drawn from a
larger ratings dataset. This complements the main MovieLens 1M analysis
by testing the methods on additional user--movie interaction data.

The main files used in this analysis are the ratings data and movie
metadata. The ratings data include user identifiers, movie identifiers,
numeric ratings, and timestamps. The movie metadata include titles and
genre labels, which are used later to interpret the resulting clusters.

The preprocessing stage transforms the raw ratings data into a form
suitable for representation learning and clustering. Because user--item
data are sparse, users with very few ratings and movies with very low
interaction frequency may be filtered out. This reduces noise and helps
ensure that the learned embeddings reflect sufficient interaction
history. Depending on the final embedding pipeline, ratings may be
retained in their original form, normalized, or converted into
implicit-feedback signals. The processed data are then used to train the
embedding model and construct user representations for clustering.

{[}WHAT TO ADD: EXACT THRESHOLDS AND FINAL SAMPLE DETAILS SHOULD BE
INSERTED ONCE THE IMPLEMENTATION IS COMPLETE.{]}

{[}REPORT THE FILTERING CRITERIA APPLIED TO USERS AND MOVIES.{]}

{[}NUMBER OF OBSERVATIONS RETAINED AFTER PREPROCESSING.{]}

{[}CONSTRUCTION OF THE SAMPLED MOVIELENS 20M BENCHMARK SUBSETS.{]}

{[}ANY NORMALIZATION OR TRANSFORMATION PROCEDURES USED IN CONSTRUCTING
THE MODEL INPUT.{]}

{[}INCLUDE A BRIEF SUMMARY OF THE SPARSITY OF THE FINAL USER--ITEM
MATRIX AND, IF RELEVANT, DESCRIPTIVE STATISTICS ABOUT RATINGS AND GENRE
FREQUENCIES.{]}

\section{Methods and Experimental
Setup}\label{methods-and-experimental-setup}

\subsection{Embeddings}\label{embeddings}

The methodology consists of two stages. First, dense latent
representations are learned from the MovieLens 1M data using an
MLP-based embedding model. Second, clustering procedures are applied to
those embeddings to determine which combinations of embedding dimension
and clustering method yield the strongest grouping structure.

The representation-learning component produces embeddings of dimensions
64, 128, and 256. These dimensions are compared to evaluate how
representation size affects clustering quality. Lower-dimensional
embeddings may encourage more compact and stable structure, while
higher-dimensional embeddings may capture subtler variation at the cost
of greater complexity.

The exact neural architecture should be documented in the final version
of the report. This subsection should specify the input representation,
the number and size of hidden layers, activation functions, optimizer,
learning rate, batch size, number of training epochs, regularization
choices, and the loss function used to train the embeddings. It should
also state explicitly that clustering is applied to user embeddings.

{[}ADD EXACT NEURAL ARCHITECTURE DETAILS.{]}

{[}ADD LOSS FUNCTION.{]}

{[}ADD OPTIMIZER, LEARNING RATE, BATCH SIZE, EPOCHS, AND
REGULARIZATION.{]}

After the embeddings are learned, clustering is performed separately in
each embedding space. The four clustering methods considered are
k-means, PD clustering, reduced k-means, and FPDC. The experimental
design applies each method to each embedding dimension in order to
compare clustering behavior across both representation size and
clustering algorithm.

The number of clusters considered, the procedure used to select that
value, the initialization strategy, stopping criteria, and the number of
repeated runs should all be stated explicitly in the final version.
These choices affect both clustering quality and computational cost.

{[}ADD NUMBER OF CLUSTERS CONSIDERED.{]}

{[}ADD CLUSTER-SELECTION PROCEDURE.{]}

{[}ADD INITIALIZATION STRATEGIES, STOPPING CRITERIA, AND NUMBER OF
REPEATED RUNS.{]}

Clustering is not performed directly on the raw sparse rating matrix.
Instead, the embedding model transforms sparse user--item behavior into
a lower-dimensional representation on which clustering is applied.

\subsection{Clustering Implementation in
R}\label{clustering-implementation-in-r}

{[}NEED CITATIONS{]}

All clustering analyses were conducted in R using standardized
user-embedding matrices derived from the MovieLens data. Embedding files
were imported with the \texttt{readr} package, and the first column
containing user identifiers was separated from the embedding
coordinates. The remaining embedding features were converted to a
numeric matrix and standardized using \texttt{scale()} so that all
dimensions contributed comparably to distance-based calculations.
Pairwise Euclidean distances were computed using \texttt{dist()} and
were used in silhouette-based evaluation. Additional data manipulation
and plotting were performed with \texttt{dplyr} and \texttt{ggplot2},
respectively.

For k-means clustering, the implementation used the base R function
\texttt{kmeans()} from the \texttt{stats} package. Candidate numbers of
clusters were evaluated over the range

\[
k \in \{2,\dots,20\}.
\]

In the tuning stage, k-means was run with \texttt{nstart\ =\ 10} and
\texttt{iter.max\ =\ 100}, and three algorithmic variants were compared:
Hartigan--Wong, Lloyd, and MacQueen. For the final fitted model, k-means
was refit with \texttt{nstart\ =\ 100} and \texttt{iter.max\ =\ 100}
using the default Hartigan--Wong algorithm. Model quality was assessed
using total within-cluster sum of squares from the fitted object,
average silhouette width computed with \texttt{silhouette()} from the
\texttt{cluster} package, and the Calinski--Harabasz index computed with
\texttt{intCriteria()} from the \texttt{clusterCrit} package. In a
separate comparison script, k-means was also run for fixed values

\[
k \in \{3,4,5,6,7,8,9,10\}
\]

using \texttt{nstart\ =\ 100}, \texttt{iter.max\ =\ 100}, and
\texttt{algorithm\ =\ "MacQueen"} to compare its partitions with reduced
k-means using the adjusted Rand index from the \texttt{mclust} package.

For reduced k-means, the implementation used the \texttt{cluspca()}
function from the \texttt{clustrd} package. Models were fit over the
grid

\[
k \in \{2,\dots,10\}, \qquad q \in \{2,\dots,6\}.
\]

Each model was fit as
\texttt{cluspca(embedding\_matrix,\ nclus\ =\ k\_val,\ ndim\ =\ q\_val,\ method\ =\ "RKM",\ rotation\ =\ "varimax")}.
For each \((k,q)\) combination, average silhouette width was computed
using \texttt{silhouette()} from the \texttt{cluster} package and the
Calinski--Harabasz index was computed with \texttt{intCriteria()} from
\texttt{clusterCrit}. The best reduced k-means model for each embedding
file was selected using the highest silhouette value, and intermediate
results were cached to \texttt{.rds} files.

Probabilistic distance clustering was implemented in R using the
\texttt{pdclust()} function on the standardized embedding matrices.
Models were fit over the range

\[
k \in \{2,3,4,5\}
\]

using Euclidean distance. For each value of \(k\), the algorithm
produced a matrix of posterior membership probabilities, from which hard
cluster assignments were obtained by assigning each observation to the
cluster with the highest probability. Cluster centers were estimated
internally by the algorithm from these probabilistic assignments.
Clustering performance was evaluated using average silhouette width
computed with \texttt{silhouette()} from the \texttt{cluster} package,
the Calinski--Harabasz index computed using \texttt{intCriteria()} from
the \texttt{clusterCrit} package, and a probability-based silhouette
measure computed with \texttt{Silh()}. All fitted models and evaluation
metrics were saved to \texttt{.rds} files to support reproducibility and
avoid recomputation.

For factorial probabilistic distance clustering (FPDC), the
implementation used the \texttt{FPDC()} function from the
\texttt{FPDclustering} package. The standardized embedding matrices were
evaluated over the grid

\[
k \in \{2,\dots,8\}, \qquad q \in \{2,\dots,5\}.
\]

Each model was fit as
\texttt{FPDC(embedding\_matrix,\ k\_val,\ 20,\ q\_val,\ nu\ =\ 10)}, so
the implementation used 20 internal iterations and a tuning parameter
fixed at

\[
\nu = 10.
\]

For each parameter combination, clustering quality was evaluated using
average silhouette width computed from hard labels with
\texttt{silhouette()} from the \texttt{cluster} package, the
Calinski--Harabasz index from \texttt{clusterCrit}, and the
probability-based silhouette returned by \texttt{Silh()} from
\texttt{FPDclustering}. The best FPDC solution for each embedding file
was selected using silhouette as the primary criterion and
Calinski--Harabasz as a secondary criterion. Each fit was cached to
\texttt{.rds} files, and summary outputs were written to \texttt{.csv}
files.

\section{Test Setup and Evaluation}\label{test-setup-and-evaluation}

Because this project is unsupervised, evaluation is based on internal
validation, stability, and interpretability rather than labeled test
outcomes. The evaluation framework therefore compares clustering
outcomes along several complementary dimensions.

The first component consists of internal validation metrics. These
metrics assess cluster compactness and separation. One key metric is the
silhouette score, which measures how similar an observation is to its
own cluster relative to the nearest competing cluster. Higher silhouette
values indicate stronger separation and better-defined structure.
Another criterion is within-cluster sum of squares, which measures
variation around assigned centroids. Smaller values indicate tighter
clustering, although this quantity generally decreases as the number of
clusters increases. Additional indices such as the Davies--Bouldin index
may also be used to broaden the comparison of compactness and
separation.

A second component concerns stability. Many clustering algorithms,
especially partition-based methods, are sensitive to initialization and
local optima. For this reason, clustering configurations should be run
multiple times, and the consistency of the resulting partitions should
be assessed. Stability can be measured using indices such as adjusted
Rand similarity. A method that yields highly variable solutions across
repeated runs may be less reliable even if its internal metrics appear
strong in a single run.

A third component is interpretability. Because the MovieLens dataset
includes movie genres and potentially user covariates, the resulting
user clusters can be examined to determine whether they correspond to
meaningful preference patterns or rating behaviors. This interpretive
step matters because a mathematically compact clustering is not
necessarily substantively useful.

The final report should state clearly which metrics were used, how many
repeated runs were performed, and how the best overall method was
selected. If dimensionality-reduction visualizations such as PCA or
t-SNE are included, they should be used as qualitative support rather
than primary evidence.

\section{Results}\label{results}

Structure:

\begin{enumerate}
\def\labelenumi{\arabic{enumi}.}
\tightlist
\item
  Effect of embedding size: compare 64 vs.~128 vs.~256.
\item
  Algorithm comparison: k-means vs.~PD vs.~reduced k-means vs.~FPDC.
\item
  Best model selection based on validation metrics.
\end{enumerate}

{[}ADD LATER: ALL NUMERIC RESULTS.{]}

{[}ADD LATER: TABLES + PLOTS.{]}

{[}ADD LATER: OBSERVATIONS (E.G., ``HIGHER DIMENSION IMPROVES SEPARATION
BUT INCREASES VARIANCE'').{]}

\section{Conclusion and Future Work}\label{conclusion-and-future-work}

This project studies clustering on learned user embeddings derived from
the MovieLens 1M dataset. By comparing embeddings of dimensions 64, 128,
and 256 across four clustering methods---k-means, PD clustering, reduced
k-means, and FPDC---the analysis evaluates how representation size and
clustering method affect unsupervised structure discovery.

The study brings together partition-based clustering, dimension
reduction, latent-factor methods, and cluster validation within a
representation-learning framework. Although recommendation performance
is not evaluated, the resulting cluster structure may still be relevant
to downstream tasks such as personalization, interpretation, and content
organization.

Future work could extend the analysis by incorporating additional
clustering methods covered in the course, such as spectral clustering,
fuzzy clustering, model-based mixture methods, or clustering with
covariates. Another extension would be to include richer user or movie
metadata in order to study clustering in augmented feature spaces. It
would also be useful to compare MLP-based embeddings with embeddings
learned through alternative architectures such as autoencoders or other
latent-variable models.

\section{Appendix}\label{appendix}

You can see our entire code at our GitHub repository.

{[}WE SHOULD CITE THE PAPERS WE READ ABOUT THE DIFFERENT CLUSTERING
METHODS!{]}

\section{References}\label{references}

\end{document}
